\ssr{ПРИЛОЖЕНИЕ А}

Презентация к научно-исследовательской работе на 8 слайдах.

\ssr{ПРИЛОЖЕНИЕ Б}

Правила  грамматики ограниченного естественного русского языка.


\small
\begin{longtable}{|p{5cm}|p{10cm}|}
	\caption{Реализуемые правила грамматики ограниченного естественного русского языка}
	\label{tab:grammar_rules}\\
	\hline
	Формальная запись правила & Описание \\
	\hline
	\endfirsthead
	
	\multicolumn{2}{p{14cm}}{Продолжение таблицы~\thetable~--- Реализуемые правила грамматики ограниченного естественного русского языка}\\
	\hline
	Формальная запись правила & Описание \\
	\hline
	\endhead
	
	\hline
	\endfoot
	
	\hline
	\endlastfoot
	
	$IP \rightarrow NP\ VP$\newline
	$IP \rightarrow VP\ NP$ &
	Группа существительного соединяется с глагольной группой образуя предикативную основу предложения\\\hline
	
	$IP \rightarrow NP\ IP$ &
	Именная группа предшествует предикативной основе предложения, дополняя ее\\\hline
	
	$IP \rightarrow NP\ ConjVP$ &
	Именная группа соединяется с глагольной группой посредством запятой или союза, образуя предикативную основу предложения\\\hline
	
	$IP \rightarrow NP\ AP$\newline
	$IP \rightarrow AP\ NP$ &
	Именная группа соединяется с группой прилагательного образуя предикативную основу без глагола в роли сказуемого\\\hline
	
	$IP \rightarrow NP\ NP$ &
	Именная группа соединяется с другой именной группой образуя предикативную основу без глагола в роли сказуемого\\\hline
	
	$IP \rightarrow PP\ NP$\newline
	$IP \rightarrow NP\ PP$ &
	Именная группа соединяется с предложной группой образуя предикативную основу без глагола в роли сказуемого\\\hline
	
	$IP \rightarrow C\ IP$ &
	Предикативная основа предложения дополняется служебным словом слева\\\hline
	
	$IP \rightarrow C\ ConjIP$ &
	Предикативная основа предложения дополняется служебным словом слева и соединяется с ним посредством союза или запятой\\\hline
	
	$IP \rightarrow AdvP\ IP$ &
	Предикативная основа предложения дополняется группой наречия слева\\\hline
	
	$IP \rightarrow IP\ ConjIP$ &
	Предикативная основа предложения соединяется с другой посредством запятой, двоеточия, точки с запятой или союза\\\hline
	
	$IP \rightarrow VP$\newline
	$IP \rightarrow AdvP$\newline
	$IP \rightarrow AP$\newline
	$IP \rightarrow NP$\newline
	$IP \rightarrow PP$ &
	Предикативная основа предложения выражена единственной синтаксической группой\\\hline
	
	$VP \rightarrow V$ &
	Глагольная группа состоит из одного глагола\\\hline
	
	$VP \rightarrow V\ NP$\newline
	$VP \rightarrow NP\ V$ &
	Глагол соединяется с именной группой, образуя глагольную группу\\\hline
	
	$VP \rightarrow V\ PP$\newline
	$VP \rightarrow PP\ V$ &
	Глагол соединяется с предложной группой, образуя глагольную группу\\\hline
	
	$VP \rightarrow VP\ PP$\newline
	$VP \rightarrow PP\ VP$ &
	Глагольная группа дополняется предложной группой\\\hline
	
	$VP \rightarrow VP\ AdvP$\newline
	$VP \rightarrow AdvP\ VP$ &
	Глагольная группа дополняется группой наречия\\\hline
	
	$VP \rightarrow V\ CP$\newline
	$VP \rightarrow CP\ V$\newline
	$VP \rightarrow V\ ConjCP$ &
	Глагол соединяется с частью предложения, вводимой союзом, либо с присоединенной частью после запятой\\\hline
	
	$VP \rightarrow VP\ ConjVP$\newline
	$VP \rightarrow VP\ VP$ &
	Две глагольные группы соединяются союзом, запятой или бессоюзной связью\\\hline
	
	$VP \rightarrow V\ NP\ PP$\newline
	$VP \rightarrow V\ NP\ NP$ &
	После глагола последовательно дополняется именной и предложной группой или двумя именными группами\\\hline
	
	$NP \rightarrow Pron$\newline
	$NP \rightarrow N$ &
	Именная группа состоит из одного местоимения или одного существительного\\\hline
	
	$NP \rightarrow AP\ N$\newline
	$NP \rightarrow N\ AP$ &
	Существительное соединяется с описательной группой, образуя именную группу\\\hline
	
	$NP \rightarrow AP\ NP$ & Именная группа дополняется описательной группой слева \\\hline
	$NP \rightarrow AP\ ConjN$ &
	Существительное соединяется с описательной группой слева через запятую или союз, образуя именную группу\\\hline
	
	$NP \rightarrow Q\ N$\newline
	$NP \rightarrow N\ Q$ &
	Числительное соединяется с существительным, образуя именную группу\\\hline
	
	$NP \rightarrow N\ N$ & Существительное соединяется с другим существительным образуя именную группу\\\hline
	$NP \rightarrow N\ NP$ & Именная группа дополняется существительным слева\\\hline
	
	$NP \rightarrow N\ PP$\newline
	$NP \rightarrow N\ CP$\newline
	$NP \rightarrow N\ ConjAP$ &
	Существительного соединяется справа с предложной группой, частью предложения, или описательной группой через союз, образуя именную группу\\\hline
	
	$NP \rightarrow NP\ ConjNP$\newline
	$NP \rightarrow NP\ NP$\newline
	$NP \rightarrow NP\ ConjCP$ &
	Именная группа дополняется другой именной группой, либо частью предложения посредством союза или запятой\\\hline
	
	$PP \rightarrow P\ NP$\newline
	$PP \rightarrow P\ Q$ &
	Предлог соединяется справа с именной группой или числительным, образуя предложную группу\\\hline
	
	$PP \rightarrow PP\ ConjPP$ &
	Предложная группа другой предложной группой посредством союза или запятой\\\hline
	
	$AP \rightarrow A$ &
	Группа прилагательного состоит из одного слова со значением признака\\\hline
	
	$AP \rightarrow AdvP\ A$\newline
	$AP \rightarrow A\ AdvP$ &
	Прилагательное соединяется с группой наречия, образуя группу прилагательного\\\hline
	
	$AP \rightarrow A\ PP$\newline
	$AP \rightarrow PP\ A$ &
	Прилагательное соединяется с предложной группой, образуя группу прилагательного\\\hline
	
	$AP \rightarrow A\ N$ &
	Прилагательное соединяется с существительным образуя группу прилагательного\\\hline
	
	$AP \rightarrow AP\ NP$& Группа прилагательного дополняется группой существительного справа\\\hline
	
	
	$AP \rightarrow AP\ ConjPP$ &
	Группа прилагательного дополняется предложной группой справа посредством союза\\\hline
	
	$AP \rightarrow AP\ ConjAP$\newline
	$AP \rightarrow AP\ AP$ &
	Группа прилагательного дополняется другой группой прилагательного посредством союза или запятой\\\hline
	
	$AdvP \rightarrow Adv$ &
	Группа наречия состоит из одного наречия\\\hline
	
	$AdvP \rightarrow AdvP\ AdvP$\newline
	$AdvP \rightarrow AdvP\ NP$\newline
	$AdvP \rightarrow AdvP\ ConjAdvP$ &
	Группа наречия дополняется другой группой наречия посредством союза или запятой или именной группой\\\hline
	
	$CP \rightarrow C\ IP$ &
	Служебное слово соединяется с предикативной основой предложения, образуя зависимую часть предложения\\\hline
	
	$ConjIP \rightarrow Comma\ IP$\newline
	$ConjIP \rightarrow Semicolon\ IP$\newline
	$ConjIP \rightarrow Colon\ IP$\newline
	$ConjIP \rightarrow C\ IP$\newline
	$ConjIP \rightarrow CommaC\ IP$ &
	Предикативная основа дополняется слева знаком препинания или служебным словом\\\hline
	
	$CommaC \rightarrow Comma\ C$ &
	Запятая соединяется с союзом образуя единую синтаксическую единицу\\\hline
	
	$ConjVP \rightarrow C\ VP$\newline
	$ConjVP \rightarrow Comma\ VP$ &
	Глагольная группа соединяется слева с союзом или запятой, образуя присоединительную глагольную группу\\\hline
	
	$ConjNP \rightarrow C\ NP$\newline
	$ConjNP \rightarrow Comma\ NP$ &
	Именная группа соединяется слева с союзом или запятой, образуя присоединительную именную группу\\\hline
	
	$ConjPP \rightarrow C\ PP$\newline
	$ConjPP \rightarrow Comma\ PP$ &
	Предложная группа соединяется слева с союзом или запятой, образуя присоединительную предложную группу\\\hline
	
	$ConjAP \rightarrow C\ AP$\newline
	$ConjAP \rightarrow Comma\ AP$ &
	Группа прилагательного соединяется слева с союзом или запятой, образуя присоединительную группу прилагательного\\\hline
	
	$ConjAdvP \rightarrow C\ AdvP$\newline
	$ConjAdvP \rightarrow Comma\ AdvP$ &
	Группа наречия соединяется слева с союзом или запятой, образуя присоединительную группу наречия\\\hline
	
	$ConjCP \rightarrow Comma\ CP$ &
	Часть предложения соединяется слева с запятой, образуя присоединительную часть предложения\\
\end{longtable}
\normalsize

\ssr{ПРИЛОЖЕНИЕ В}

Правила согласования применяемые в ограниченном русском языке.

\small
\begin{longtable}{|p{4.5cm}|p{9.5cm}|}
	\caption{Правила согласования ограниченного естественного русского языка}
	\label{tab:agreement_rules}\\
	\hline
	Формальная запись правила согласования & Описание \\
	\hline
	\endfirsthead
	
	\multicolumn{2}{p{14cm}}{Продолжение таблицы~\thetable~--- Правила согласования ограниченного естественного русского языка}\\
	\hline
	Формальная запись правила согласования & Описание \\
	\hline
	\endhead
	
	\hline
	\endfoot
	
	\hline
	\endlastfoot
	
	$IP:\ NP \leftrightarrow VP$\newline $IP:\ VP \leftrightarrow NP$\newline $IP:\ NP \leftrightarrow ConjVP$  &
	Именная и глагольная группы согласуются в числе. Дополнительно в прошедшем времени согласование в роде, в настоящем времени~---~в лице. Если глагольная группа предшествует именной и выражена неопределенной формой, именная группа рассматривается как дополнение\\\hline
	
	$IP:\ NP \leftrightarrow AP$\newline
	$IP:\ AP \leftrightarrow NP$ &
	В предложении без глагола в роли сказуемого именная группа и группа прилагательного согласуются в роде, числе и падеже\\\hline
	
	$IP:\ NP \leftrightarrow NP$ &
	В предложении без глагола в роли сказуемого две именные группы согласуются в числе и роде\\\hline
	
	$IP:\ NP \leftrightarrow PP$\newline
	$IP:\ PP \leftrightarrow NP$ &
	В сочетании именной и предложной группы именная группа принимает форму именительного падежа\\\hline
	
	$IP:\ NP_{nomn} \leftrightarrow IP$ &
	Именная группа в именительном падеже перед предикативной основой согласуется с ней в числе; в прошедшем времени~---~в роде; в настоящем времени~---~в лице\\\hline
	
	$NP:\ AP \leftrightarrow N$\newline
	$NP:\ AP \leftrightarrow NP$\newline
	$NP:\ AP \leftrightarrow ConjN$\newline
	$NP:\ N \leftrightarrow AP$ &
	Группа прилагательного согласуется с именной группой или присоединительной именной группой, или существительным в роде, числе и падеже\\\hline
	
	$NP:\ N \rightarrow N_{gent}$\newline
	$NP:\ N \rightarrow NP_{gent}$ &
	Существительное или именная группа после существительного принимают форму родительного падежа\\\hline
	
	$VP:\ V \rightarrow NP$ &
	Именная группа после глагола не принимает форму именительного падежа. После переходных глаголов именная группа принимает форму только винительного падежа\\\hline
	
	$PP:\ P \rightarrow NP$ &
	Падеж именной группы после предлога определяется видом предлога на основе словаря\\\hline
	
	$NP:\ NP \leftrightarrow NP$\newline
	$NP:\ NP \leftrightarrow ConjNP$ &
	Две именные группы согласуются в падеже\\\hline
	
	$VP:\ VP \leftrightarrow VP$\newline
	$VP:\ VP \leftrightarrow ConjVP$ &
	Две глагольные группы согласуются во времени, числе и лице\\\hline
	
	$AP:\ AP \leftrightarrow AP$\newline
	$AP:\ AP \leftrightarrow ConjAP$ &
	Две группы прилагательных согласуются в роде, числе и падеже\\\hline
	
	$IP:\ VP \not\rightarrow infn$ &
	Неопределенная форма глагола не может быть использована как предикативная основа \\\hline
	
	$NP:\ Q \leftrightarrow N$\newline
	$NP:\ N \leftrightarrow Q$ &
	Числительное и существительное согласуются в роде числа\\\hline
	
	$NP:\ Q \leftrightarrow AP \leftrightarrow N$ &
	В именной группе с числительным группа прилагательного согласуется с ним в числе и падеже\\\hline
	
\end{longtable}
\normalsize